% Indicate the main file. Must go at the beginning of the file.
% !TEX root = ../main.tex

%----------------------------------------------------------------------------------------
% CHAPTER TEMPLATE
%----------------------------------------------------------------------------------------


\chapter*{Overview of Applied AI Tools}
\addcontentsline{toc}{chapter}{Overview of Applied AI Tools}

\begin{itemize}

    \item \textbf{Anthropic (2025).} \textit{Claude Pro (Sonnet 4.5).}\\
    Available at: \url{https://claude.ai/} \\
    Applications: Coding assistance and programming support, including code structure optimization, debugging, algorithm explanation, and enhancing understanding of best practices. Also used for text-related tasks including generating, editing, and refining academic, professional, and creative content.

    \item \textbf{DeepL (2025).} \textit{DeepL Translator and DeepL Write.}\\
    Available at: \url{https://www.deepl.com/translator} \\
    Applications: Text translation, grammar correction, synonym search, translation between German and English, and suggestions for academic writing.

    \item \textbf{Grammarly (2025).} \textit{Grammarly Premium.}\\
    Available at: \url{https://www.grammarly.com/} \\
    Applications: Grammar correction, sentence restructuring, and suggestions to improve clarity and tone.

    \item \textbf{Microsoft (2025).} \textit{Microsoft Copilot.}\\
    Available at: \url{https://www.microsoft.com/en-us/microsoft-copilot} \\
    Applications: Coding assistance and debugging support for Python and general programming tasks. Assisted with the clarification of scientific and medical topics and provided LaTeX formatting support.

    \item \textbf{NotebookLM (2025).} \textit{NotebookLM.}\\
    Available at: \url{https://notebooklm.google.com/} \\
    Applications: Searching and summarizing relevant information contained within uploaded documents (e.g., research papers, lecture notes). Used to extract key ideas, clarify theoretical concepts, and create structured summaries. Supports understanding of domain-specific content by providing contextual explanations based on user-supplied sources.

    \item \textbf{Perplexity (2025).} \textit{Perplexity AI.}\\
    Available at: \url{https://www.perplexity.ai/} \\
    Applications: Assisting with research by retrieving and summarizing information from diverse sources, including academic and scientific materials. Used to explore unfamiliar topics, gather context for literature understanding, and obtain concise explanations of technical concepts. Supports guided question--answer workflows for efficient knowledge acquisition.

\end{itemize}
